\color{unidarkgreen}{\subsection{ Устройство резервирования отказа выключателя}}
\color{black}

\begin{enumerate}

\item
Устройство резервирования отказа выключателя (УРОВ) предназначено для выполнения функции ближнего резервирования защит. В случае отказа выключателя повреждённого присоединения УРОВ обеспечивает отключение смежных выключателей с целью селективного отделения повреждённого участка электрической сети.

\item
Для выявления факта отказа выключателя предусмотрен контроль наличия тока через выключатель. Контроль осуществляется измерительным органом УРОВ, порог срабатывания которого задается уставкой \textbf{«Iср УРОВ»}.

\item
Входной сигнал ФСУ \textit{«УРОВ В1(ЛВ) ВН/УРОВ: Пуск УРОВ»} предусмотрен для пуска УРОВ от защит, срабатывание которых не сопровождается увеличением тока или сопровождается протеканием малых токов КЗ, недостаточных для срабатывания токовых ИО УРОВ.

\item
Действие УРОВ на себя осуществляется от сигнала внешнего пуска \textit{«УРОВ В1(ЛВ) ВН/УРОВ: Пуск УРОВ»} с подхватом от ИО УРОВ.
Сигнал срабатывания УРОВ на себя \textit{«УРОВ В1(ЛВ) ВН/УРОВ: Срабатывание на себя»} формируется при устанавленном программном ключе \textbf{«УРОВ Действ на себя»} в положение \textit{«Введено»} с выдержкой времени \textbf{«Тср повт УРОВ»}.

\item
Действие УРОВ на смежные выключатели осуществляется от сигнала внешнего пуска \textit{«УРОВ В1(ЛВ) ВН/УРОВ: Пуск УРОВ»} с подхватом от ИО УРОВ. Предусматривается возможность формирования выходного действия УРОВ на смежные выключатели с контролем сигнала о включенном положении выключателя \textit{«ФПВ: В включен»}.

\item
Ввод в работу контроля положения выключателя осуществляется устанавкой программного ключа \textbf{«УРОВ конт. выкл.»} в положение \textit{«Введено»}.

\item
Действие УРОВ на смежные выключатели выполняется с выдержкой времени \textbf{«Тср УРОВ»}.

\item
Функция УРОВ может быть оперативно выведена из работы от внешнего сигнала \textit{«УРОВ В1(ЛВ) ВН/УРОВ: Вывод»}, назначенного на соответствующий дискретный вход или функциональную клавишу, или с помощью программного ключа \textbf{«УРОВ»}, установленного в положении \textit{«Выведено»}.

\item
Функциональная схема логики УРОВ приведена на рисунке \ref{pict_urov:UROV}. Параметры настройки функции УРОВ представлены в таблице \ref{app_set:urov}.

\medskip
\begin{figure}[H]
\centering
\includegraphics[width=0.7\textwidth,height=0.7\textheight,keepaspectratio]{\fbpath/10.02 АУВ/УРОВ/img/UROV.pdf}
\captionof{figure}{Функциональная схема логики УРОВ}\label{pict_urov:UROV}
\end{figure}

\end{enumerate}